% Generated draft: requires mathematical and source review.
\documentclass[11pt]{article}
\usepackage[margin=25mm]{geometry}
\usepackage{fontspec}
\setmainfont{Latin Modern Roman}
\usepackage{amsmath,amssymb,mathtools,unicode-math}
\setmathfont{Latin Modern Math}
\usepackage{enumitem,needspace,xurl,hyperref}
\setlength{\parindent}{0pt}
\setlength{\parskip}{5pt}
\setlength{\emergencystretch}{3em}
\newcommand{\N}{\mathbb{N}}
\newcommand{\Z}{\mathbb{Z}}
\newcommand{\Q}{\mathbb{Q}}
\newcommand{\R}{\mathbb{R}}
\newcommand{\C}{\mathbb{C}}
\newcommand{\F}{\mathbb{F}}
\newcommand{\A}{\mathbb{A}}
\newcommand{\PP}{\mathbb{P}}
\newcommand{\E}{\mathbb{E}}
\newcommand{\eps}{\varepsilon}
\newcommand{\dd}{\,\mathrm{d}}
\begin{document}
\section*{UnsolvedMath ID: 2307078}
\textbf{Source problem code:} catalogue entry
\subsection{Definitions and mathematical statement}
Let $(z_n)_{n\geq1}$ be pairwise distinct nonzero complex numbers with
\[
 \sum_{n=1}^{\infty}\frac1{|z_n|}<\infty.
\]
Then $|z_n|\to\infty$, and the series
\[
 F(z)=\sum_{n=1}^{\infty}\frac1{z-z_n}
\]
converges locally uniformly away from the $z_n$, defining a meromorphic function with a simple pole of residue one at each $z_n$.

L. A. Rubel's existence question is whether such a sequence can make $F$ zero-free:
\[
 \exists(z_n)_{n\geq1}\text{ as above}\ \forall z\in\C\setminus\{z_n:n\geq1\},
 \qquad F(z)\ne0.
\]
The source says ``for all $z\in\C$,'' although the displayed series is undefined at its poles. The precise interpretation above either excludes the poles or equivalently views $F$ meromorphically, with poles not counted as zeros; this domain qualification must be confirmed rather than hidden.

The source interprets $F$ as the field of unit-charged wires perpendicular to the plane and asks whether an equilibrium must always exist. Its formal logarithmic potential $\sum_n\log|z-z_n|$ needs additive renormalization: those summands do not tend to zero. A well-defined normalized potential away from the charges is
\[
 U(z)=\sum_{n=1}^{\infty}\log|1-z/z_n|,
 \qquad \partial_xU-i\partial_yU=F(z).
\]
Thus a field equilibrium away from charges is exactly a zero of $F$, with the complex-gradient convention explicit. This normalization explains the physical statement without asserting convergence of the raw potential.

The source also proposes an analogous question in $\R^n$ for $n\geq3$, but supplies no corresponding potential, convergence condition or precise higher-dimensional formulation; that suggested extension remains unspecified.
\subsection{Short English statement}
Can a summably sparse infinite set of equal plane charges produce a field with no equilibrium away from the charges? The field series is well defined under the stated reciprocal-distance sum; the original all-plane wording needs a pole convention and the raw logarithmic potential needs normalization. The supplied triage reports the question open in Hayman's 2018 update and no definitive resolution located; no new example or current-status assessment is claimed.
\subsection{Sources}
\begin{itemize}
\item UnsolvedMath permanent ID 2307078, source code catalogue entry, ``Research Problems in Function Theory --- Problem 7.78.'' Registry curation and metadata: Ulam AI and UnsolvedMath Contributors, CC BY 4.0; updated 2026-08-06.
\item Hayman, \emph{Research Problems in Function Theory} (2018 source update), Problem 7.78, attributed to L. A. Rubel: \url{https://arxiv.org/abs/1809.07200}. Supplied source-TeX extraction accessed 2026-07-29. The metadata retains the update for status triage and marks status/release and rights review NEEDS\_REVIEW.
\item No additional authoritative result or reference is supplied for a resolution or the higher-dimensional suggestion. No missing citation is invented.
\end{itemize}
\end{document}
